\section{Experimental Results}
\label{sec:experiments}
In this section, we validate our implementations and benchmark some RL, tree search, and model-based game solving algorithms on our implemented environments. We first validate the ability of our JAX-based implementation of CFR to find near-optimal strategies for solved games, as well as standard RL, RL augmented with NFSP, and Deep CFR algorithms. We then compare RL+NFSP with ReBeL baselines on Texas Hold'Em Poker, a well-studied medium for game-theoretic algorithms. Finally, we benchmark our remaining environments with RL+NFSP baselines.

\begin{figure}[htbp]
    \centering
    \begin{subfigure}[b]{0.45\textwidth}
        \centering
        \includegraphics[width=\textwidth]{02_images/baselines/kuhn.pdf}
        \caption{Kuhn Poker}
        \label{fig:kuhn}
    \end{subfigure}
    % \hfill
    \begin{subfigure}[b]{0.45\textwidth}
        \centering
        \includegraphics[width=\textwidth]{02_images/baselines/leduc.pdf}
        \caption{Leduc Poker}
        \label{fig:leduc}
    \end{subfigure}
    \caption{Exploitability of Kuhn and Leduc Poker. Each algorithm is trained in self-play with 1e7 samples and evaluated across 10 seeds.}
    \label{fig:kuhn_leduc}
\end{figure}

Figure~\ref{fig:kuhn_leduc} shows the exploitability of various algorithms on Kuhn and Leduc Poker. The exploitability of a strategy profile is defined as net utility gained by an optimal policy playing against that strategy profile. Lower exploitability is better; an exploitability of $0$ means a strategy profile is a Nash equilibrium~\citep{nash1950equilibrium}. We note the exploitability of a uniform random policy on Kuhn ($0.4583$) and Leduc ($2.3736$), which matches values from \citet{kawamura2017neural} and OpenSpiel. We also include the exploitability obtained by $100$ iterations of CFR. Deep CFR achieves near-convergence on Kuhn Poker with $100$ iterations, but struggles with Leduc Poker. This may be because truncating the relatively shallow game tree of Leduc Poker may exaggerate errors from function approximation.

We benchmark the exploitability of two RL algorithms, Proximal Policy Optimization (PPO)~\citep{schulman2017proximal} and Parallelized Q-Network (PQN)~\citep{gallici2024simplifying}, on these environments. We choose PPO and PQN as they are simple to implement, well-known by the RL research community, and are well-suited to parallelized environment training and GPU acceleration. While PPO and PQN are able to achieve decent results in Kuhn Poker (beating Uniform Random by a fair margin), they tend to exhibit oscillatory behavior. Cycling is a well-studied phenomenon in two-player zero-sum games~\citep{shapley1963some, balduzzi2018mechanics} which is worsened in self-play settings due to bootstrapping with self-produced training data~\citep{heinrich2016deep}.

We also implement PPO-NFSP and PQN-NFSP, which are versions of NFSP with PPO and PQN as the best response learners respectively alongside a supervised anticipatory network. We note that they achieve better results than vanilla PPO and PQN in both environments, but still fall short of optimality in Leduc Poker. 

\begin{figure}[htbp]
    \centering
    \begin{subfigure}[b]{0.4\textwidth}
        \centering
        \includegraphics[width=\textwidth]{02_images/baselines/hul.pdf}
        \caption{Heads-Up Texas Hold'Em Limit}
        \label{fig:hul}
    \end{subfigure}
    % \hfill
    \begin{subfigure}[b]{0.4\textwidth}
        \centering
        \includegraphics[width=\textwidth]{02_images/baselines/hunl.pdf}
        \caption{Heads-Up Texas Hold'Em No-Limit}
        \label{fig:hunl}
    \end{subfigure}
    \caption{Final performance of PPO-NFSP, PQN-NFSP, and ReBeL against pre-trained ReBeL evaluation opponent in Heads-Up Texas Hold'Em. Each algorithm is trained in self-play with 1e7 samples and evaluated across 10 seeds. Evaluation opponent is trained in self-play with 5e6 samples. Dashed lines indicate break-even performance.}
    \label{fig:holdem}
\end{figure}

In Figure~\ref{fig:holdem} we test the performance of baseline algorithms in Heads-Up (two-player) Texas Hold'Em Poker. In order to provide standardized performance results, we use ReBeL trained in self-play for 5e6 steps as an evaluation opponent. As expected, ReBeL performs well when tested against itself, due to access to more training samples and exploitation strategies learned during self-play. PPO-NFSP and PQN-NFSP are typically unable to reach break-even performance despite training on double the number of samples as the ReBeL baseline. This is likely due to ReBeL's well-tuned formulation for Hold'Em Poker, which utilizes both depth-limited CFR for subgame solving and function approximation, thereby leveraging game knowledge which is unavailable in model-free RL.

\begin{figure}[htbp]
    \centering
    \begin{subfigure}[b]{0.3\textwidth}
        \centering
        \includegraphics[width=\textwidth]{02_images/baselines/five_card_draw.pdf}
        \caption{Five Card Draw}
        \label{fig:5cd}
    \end{subfigure}
    \hfill
    \begin{subfigure}[b]{0.3\textwidth}
        \centering
        \includegraphics[width=\textwidth]{02_images/baselines/seven_card_stud.pdf}
        \caption{Seven Card Stud}
        \label{fig:7cs}
    \end{subfigure}
    \hfill
    \begin{subfigure}[b]{0.3\textwidth}
        \centering
        \includegraphics[width=\textwidth]{02_images/baselines/werewolf.pdf}
        \caption{Werewolf}
        \label{fig:werewolf}
    \end{subfigure}

    
    
    \begin{subfigure}[b]{0.3\textwidth}
        \centering
        \includegraphics[width=\textwidth]{02_images/baselines/goofspiel.pdf}
        \caption{Goofspiel}
        \label{fig:goofspiel}
    \end{subfigure}
    \hfill
    \begin{subfigure}[b]{0.3\textwidth}
        \centering
        \includegraphics[width=\textwidth]{02_images/baselines/kemps.pdf}
        \caption{Kemps}
        \label{fig:kemps}
    \end{subfigure}
    \hfill
    \begin{subfigure}[b]{0.3\textwidth}
        \centering
        \includegraphics[width=\textwidth]{02_images/baselines/bluff.pdf}
        \caption{Bluff}
        \label{fig:bluff}
    \end{subfigure}
    \caption{Final performance of PPO-NFSP and PQN-NFSP against pre-trained PPO-NFSP evaluation opponent in BluffJAX games. Each algorithm is trained in self-play with 1e7 samples and evaluated across 10 seeds. Evaluation opponent is trained in self-play with 5e6 samples. Dashed lines indicate break-even performance.}
    \label{fig:other_games}
\end{figure}

Finally, in Figure~\ref{fig:other_games} we evaluate the performance of baseline algorithms in the remaining 6 games in BluffJAX. Note that we do not test ReBeL here as it utilizes model-based subgame solving, which is currently only available for Heads-Up Hold'Em variants. Therefore, we use PPO-NFSP trained in self-play for 5e6 steps as an evaluation opponent. Across the 6 games tested, PPO-NFSP had higher average performance in 3 with higher margins relative to PQN. Notably, the results for Five Card Draw, Seven Card Stud, and Werewolf indicate that the baselines were able to gain a noticeable advantage in certain games despite their lack of game knowledge exploitation. However, both algorithms struggled to gain an advantage in Goofspiel, likely because performing well requires a stochastic policy. Finally, both algorithms failed to break even in Kemps and Bluff. It should be noted that, by virtue of their game mechanics, these are the two environments in BluffJAX which do not enforce a game-ending condition in finite time (i.e. games can end via reaching a horizon). Thus, low win rates in these games are due to draws in horizon-timeouts. Kemps introduces difficult mechanics: not only must agents learn to signal when they have built a hand, but their partner must also correctly register that signal - certainly a strong challenge in protocol formation for RL. Similarly, Bluff requires introduces a harsh `theory of mind' challenge to adversarial RL, as agents must successfully predict when opponents are lying.